problem:  
Let \((M^4,g)\) be a closed oriented Riemannian 4‑manifold.  
Denote by \(P_g\) the Paneitz operator and \(Q_g\) the Branson \(Q\)‑curvature. For a smooth function \(f\in C^\infty(M)\) satisfying  
\[
\int_M f\, dV_g = \int_M Q_g \, dV_g = 8\pi^2\,k,\quad k\in\mathbb{Z},
\]  
consider the prescribed \(Q\)-curvature equation  
\[
P_g u + 2 Q_g = 2 f e^{4u},\qquad u\in C^\infty(M).
\]

Assume \(P_g\) has **trivial kernel** and possibly changes sign (i.e. its first nonzero eigenvalue can be positive or negative).

**Problem:**  
Define the **local solvability radius**
\[
\lambda(M,[g]) := \sup\Bigl\{\lambda>0 : \text{the equation has a smooth solution for all } f \text{ with } \|f-Q_g\|_{L^2(M,g)} < \lambda\Bigr\}.
\]
Investigate the analytical‑geometric structure and possible universal lower bounds of \(\lambda(M,[g])\):

1. **Conformal invariance and spectral dependence:**  
   Does \(\lambda(M,[g])\) depend only on the conformal class \([g]\)? Can it be expressed (or bounded) in terms of conformally invariant spectral quantities—such as the first nonzero eigenvalue of \(P_g\) or geometric integrals involving its Green’s function?

2. **Singular threshold behavior:**  
   If a sequence of conformal classes \([g_i]\) satisfies \(\lambda(M,[g_i]) \to 0\), must this degeneration correspond to a geometric or topological phenomenon—such as concentration of the Green’s function of \(P_{g_i}\) or blow‑up of the Weyl curvature \(W_{g_i}\)?  

3. **Bifurcation and conformal degeneration:**  
   Is there a bifurcation in the moduli space of conformal structures when the sign of the first eigenvalue of \(P_g\) changes—causing a discontinuous drop of \(\lambda(M,[g])\)—analogous to the transition in Yamabe metrics when the scalar curvature sign changes?  

Why is it a "good" problem:  
This revised formulation corrects prior inconsistencies and establishes a precise, analytically sound research question that remains open. It introduces the notion of a *local solvability radius* \(\lambda(M,[g])\), a quantitative measure of stability for the prescribed \(Q\)‑curvature equation. The problem is “good” because:

- It defines a new conformally invariant analytic quantity that could unify nonlinear PDE behavior with spectral geometry in dimension four.  
- It asks for a fundamental relationship between conformal invariance, spectral data (Paneitz operator), and curvature analysis, aiming to characterize when small perturbations of \(Q_g\) preserve solvability.  
- It isolates the conjectural threshold where non‑coercivity or geometric degeneration leads to failure of solvability, mirroring boundary phenomena in Yamabe theory but for a fourth‑order conformal setting.  
- The formulation is simple—quantifying how far one can perturb \(Q_g\)—yet it probes deep analytic, geometric, and topological structures in conformal 4‑geometry, potentially leading to a perturbative stability theory for \(Q\)‑curvature.